% Generated draft: requires mathematical and source review.
\documentclass[11pt]{article}
\usepackage[margin=25mm]{geometry}
\usepackage{fontspec}
\setmainfont{Latin Modern Roman}
\usepackage{amsmath,amssymb,mathtools,unicode-math}
\setmathfont{Latin Modern Math}
\usepackage{enumitem,needspace,xurl,hyperref}
\setlength{\parindent}{0pt}
\setlength{\parskip}{5pt}
\setlength{\emergencystretch}{3em}
\newcommand{\N}{\mathbb{N}}
\newcommand{\Z}{\mathbb{Z}}
\newcommand{\Q}{\mathbb{Q}}
\newcommand{\R}{\mathbb{R}}
\newcommand{\C}{\mathbb{C}}
\newcommand{\F}{\mathbb{F}}
\newcommand{\A}{\mathbb{A}}
\newcommand{\PP}{\mathbb{P}}
\newcommand{\E}{\mathbb{E}}
\newcommand{\eps}{\varepsilon}
\newcommand{\dd}{\,\mathrm{d}}
\begin{document}
\section*{UnsolvedMath ID: 5100067}
\textbf{Source problem code:} catalogue entry
\subsection{Definitions and mathematical statement}
Fix a nondegenerate billiard ellipse $E$ with center $O$ and foci $f_1,f_2$, and a fixed confocal caustic $C$ admitting a Poncelet family of $N$-periodic billiard polygons. Let $P=(P_1,\ldots,P_N)$ range over this family, with indices modulo $N$: its vertices lie on $E$, its sides are tangent to $C$, and the billiard reflection law holds at every vertex. Let $P'$ be the polygon of consecutive intersections of tangent lines to $E$ at the vertices of $P$, and $P''$ the polygon of caustic tangency points.

For $j=1,2$, unit-circle inversion centered at $f_j$ is
\[
I_j(x)=f_j+\frac{x-f_j}{\|x-f_j\|^2}\qquad(x\ne f_j),
\]
and $P_j^\dagger$ is the polygon with vertices $I_j(P_i)$. The source denotes certain focal dual areas by $A_{j,\mathrm{dual}}^\dagger$. It does not define the focal dual construction or specify the order of duality and inversion; that notation is retained without inventing a construction.

For every even $N$ for which the stated family and focal dual objects are defined, prove invariant $k_{907,b}$:
\[
\frac{A_{1,\mathrm{dual}}^\dagger(P)}{A_{2,\mathrm{dual}}^\dagger(P)}=1
\qquad\text{for every }P\text{ in the fixed-caustic family},
\]
where the ratio requires $A_{2,\mathrm{dual}}^\dagger(P)\ne0$. Equivalently, the two areas are equal whenever both are defined and the denominator is nonzero. The constant is exactly one, rather than merely an unspecified constant depending on the family. The source gives no signed-area convention or treatment of degenerate members, and the missing dual definition prevents a fully geometric formulation of these area functions.
\subsection{Short English statement}
For even-period elliptic-billiard families with a fixed confocal caustic, prove that the two focal dual inverted areas have ratio exactly one. The supplied extract does not define the dual construction or its area conventions.
\subsection{Sources}
\begin{itemize}
\item Dan Reznik, Ronaldo Garcia, and Jair Koiller, \emph{Eighty New Invariants of N-Periodics in the Elliptic Billiard}, 2021 source-list version, table row $k_{907,b}$ in ``Pairs of inversive objects,'' \url{https://arxiv.org/abs/2004.12497}; accessed 2026-07-29, source-TeX extraction. The supplied summary says the proven column is ``?'' and the observation was experimentally detected in April--May 2020.
\item Follow-up references mentioned in the supplied summary: Roitman--Garcia--Reznik, bicentric paper, arXiv:2103.11260; Reznik--Garcia--Helman, inversive-triangle paper, arXiv:2012.03020; Garcia--Reznik, self-intersected paper, arXiv:2011.06640. It also begins a Garcia--Koiller--Reznik ``New\ldots'' reference, truncated before identifying details.
\item Local registry \texttt{data/unsolvedmath/problems.json}, permanent UnsolvedMath ID $5100067$, source code catalogue entry, \emph{Elliptic-billiard invariant $k_{907,b}$}. Dataset attribution: Ulam AI and the UnsolvedMath Contributors.
\end{itemize}
\end{document}
